% Les quatre parcours de lecture, en regard de la section « Parcours de
% lecture » de FrontBackmatter/CommentLire.tex. Meme palette que carte.tex:
% chaque couloir est une suite de jalons, du point d'entree au chapitre vise.
\definecolor{lctx}{RGB}{90,90,90}
\definecolor{lcinq}{RGB}{40,70,140}
\definecolor{lsix}{RGB}{50,110,120}
\definecolor{lsept}{RGB}{125,125,125}
\begin{tikzpicture}[
    >={Stealth[length=4pt,width=3pt]},
    line join=round,
    font=\footnotesize,
    lane/.style={anchor=east, align=right, text width=2.3cm,
                 font=\scriptsize\bfseries, inner sep=0pt},
    jalon/.style={draw, rounded corners=2pt, align=center, font=\scriptsize,
                 text width=1.75cm, inner sep=4pt, minimum height=0.95cm},
    goal/.style={jalon, very thick},
    hop/.style={thin, ->, shorten >=1pt, shorten <=1pt},
    note/.style={font=\tiny, text=lctx, align=center},
  ]

  % abscisses des quatre colonnes de jalons
  \def\ca{3.45} \def\cb{6.05} \def\cc{8.65} \def\cd{11.25}

  %--------------------------------------------------- 1. radar / sonar
  \node[lane, text=lcinq] at (2.30,0) {Traitement du signal RADAR ou SONAR};
  \node[jalon, draw=lcinq!55, text=lcinq] (a1) at (\ca,0) {bilans des ch.~\textsc{ii}.1 et \textsc{ii}.2};
  \node[jalon, draw=lcinq!55, text=lcinq] (a2) at (\cb,0) {\S3.2\\cône SPD};
  \node[jalon, draw=lcinq!55, text=lcinq] (a3) at (\cc,0) {\S3.3\\moyennes};
  \node[goal, draw=lcinq, text=lcinq, fill=lcinq!7] (a4) at (\cd,0) {ch.~\textsc{iii}.1\\détection};
  \draw[hop, lcinq] (a1) -- (a2);
  \draw[hop, lcinq] (a2) -- (a3);
  \draw[hop, lcinq] (a3) -- (a4);
  \node[note, anchor=north] at ($(a3)!0.5!(a4)+(0,-0.50)$)
    {\S4.1 en appoint\\(tests de Rao)};

  %--------------------------------------------------- 2. geometres
  \node[lane, text=lsix] at (2.30,-2.05) {Géométrie riemannienne, apprentissage sur variétés};
  \node[jalon, draw=lsix!55, text=lsix] (b1) at (\ca,-2.05) {\S2.1.2\\elliptiques};
  \node[jalon, draw=lsix!55, text=lsix] (b2) at (\cb,-2.05) {\S2.1.3\\M-estimation};
  \node[jalon, draw=lsix!55, text=lsix] (b3) at (\cc,-2.05) {ch.~\textsc{ii}.4\\Fisher};
  \node[goal, draw=lsix, text=lsix, fill=lsix!7] (b4) at (\cd,-2.05) {ch.~\textsc{iii}.2 et \textsc{iii}.3};
  \draw[hop, lsix] (b1) -- (b2);
  \draw[hop, lsix] (b2) -- (b3);
  \draw[hop, lsix] (b3) -- (b4);

  %--------------------------------------------------- 3. lecture rapide
  \node[lane, text=lsept] at (2.30,-4.10) {Lecture rapide};
  \node[jalon, draw=lsept!60, text=lsept] (c1) at (\ca,-4.10) {curriculum vitae};
  \node[jalon, draw=lsept!60, text=lsept] (c2) at (\cb,-4.10) {bilans des ch.~\textsc{ii}.1 à \textsc{ii}.4};
  \node[goal, draw=lsept, text=lsept, fill=lsept!7] (c3) at (\cc,-4.10)
    {intro.\ et bilans de la partie \textsc{iii}};
  \draw[hop, lsept] (c1) -- (c2);
  \draw[hop, lsept] (c2) -- (c3);
  \node[note, anchor=west] at ($(c3.east)+(0.35,0)$)
    {les blocs \emph{Contributions}\\donnent les publications};

  %--------------------------------------------------- 4. lecture lineaire
  \node[lane, text=lctx] at (2.30,-5.85) {Lecture linéaire};
  \node[jalon, draw=lctx!55, text=lctx, text width=8.65cm, minimum height=0.7cm,
        anchor=west] (d1) at (\ca-0.95,-5.85)
    {l'intégralité, dans l'ordre du mémoire: chaque chapitre de l'état des lieux
     naît d'une insuffisance du précédent};

  %--------------------------------------------------- reperes de colonne
  \node[note, anchor=south] at (\ca,0.75) {point d'entrée};
  \node[note, anchor=south] at (\cd,0.75) {chapitre visé};
  \draw[lctx!30, dashed, thin] (\ca,0.65) -- (\ca,-5.35);
  \draw[lctx!30, dashed, thin] (\cd,0.65) -- (\cd,-2.75);
\end{tikzpicture}
